\documentclass[../Vibro-Acoustic Signatures and Detection Algorithms for Stored Product Insects.tex]{subfiles}
\begin{document}

\chapter{Introduction}

Stored products, such as grains, legumes, spices, and processed foods represent a significant portion of global food supply and trade. According to the Food and Agriculture Organization (FAO) of the United Nations, 3029 million tonnes of cereal products, including rice and wheat, were produced globally in 2025, of which 880 million tonnes were traded internationally \cite{FoodOutlookBiannual2025}. These products are susceptible to infestation by insects, most commonly referred to as Stored Product Pests (SPPs). An undetected or unmitigated insect proliferation can foster quantitative and qualitative damage to the commodity as the pests consume and contaminate the material, impacting its supply, salability, and presenting food safety concerns. Furthermore, the introduction of non-native species into ecological environments can exacerbate negative effects and in the past have led to financial expenses from losses, eradication, and prevention estimated at \$26 billion a year since 2010 for North America alone \cite{crystal-ornelasEconomicCostsBiological2021}. Efficaciously inspecting products for pests throughout the global supply chain is crucial to ensure the integrity of the product, deterring greater infestation, and securing regional agricultural and environmental ecologies.

A strategic gateway to control the entry of SPPs into the United States is by inspecting vulnerable materials at the 328 Ports of Entry (PoE) including seaports, airports, and land border crossings that annually oversee the border traffic of 11 million maritime containers, 11 million truck shipments, and 2.7 million rail shipments, respectively \cite{PortsEntryUS2025}. Currently, this responsibility is entrusted to the United States Customs and Border Protection (CBP) Agriculture Specialists (CBPAS) whose main method of stored product inspection relies on manual visual examination. This process is labor-intensive, time-consuming, and prone to human error as individuals sift through large quantities of product looking for signs of at least one live or dead insect, fragments, or other indicators of infestation. The inspection is further complicated by the fact that many insects are small, cryptic, and may be concealed within the kernels of grains and seeds \cite{service2015a}. 

Various alternative methods have been explored to enhance and automate the inspection process and detection capabilities. Physical methods, use probes or counters to capture and enumerate the insects as they pass through a receptacle but are limited to detecting only mobile insects or can damage the product in the process. Chemical sensors monitor volatile organic compounds (VOCs) including oxygen, carbon dioxide, amino acid, uric acid, and pheromones that are emitted by the insect but the sensors are expensive and require frequent calibration. Spectral methods, such as X-ray imaging and near-infrared spectroscopy (NIRS), can non-destructively scan the product to identify anomalies that may indicate the presence of insects but are also costly, complex, and require specialized training \cite{muller-blenkleMethodAcousticStorage2023}.

Passive acoustics has been a reliable and relatively low-cost alternative to even the most advanced sensing technologies to provide detection, localization, tracking, and classification of subjects including wood boring insects \cite{flynnAcousticMethodsInvasive2016, sutinVibroacousticMethodsInsect2017, sutinAutomatedAcousticDetection2019}, marine mammals, scuba divers and surface vessels \cite{sutinStevensPassiveAcoustic2010, sutinStevensPassiveAcoustic2013, pollaraPassiveAcousticMethods2017, sedunovLowSizeCostAcoustic2022, salloumLowcostMultimodalIntegrated2022, sedunovStevensPassiveAcoustic2022a}, low-flying aircraft \cite{sutinAcousticDetectionTracking2013, sedunovPassiveAcousticLocalization2013}, and unmanned aerial vehicles (UAV) \cite{yakubovskiyFeatureExtractionAcoustic2015, salloumAcousticSystemLow2015, sedunovUAVPassiveAcoustic2018, sedunovStevensDroneDetection2019,kadyrovImprovementsStevensDrone2022}. 

Acoustic and vibration sensors, such as traditional microphones, piezoelectric discs, accelerometers, ultrasonic transducers, micro-electromechanical systems (MEMS), and laser vibrometers, have been effectively used to capture the sound waves and vibrations generated by insects as they move, communicate, masticate, and procreate in the material \cite{mankinPerspectivePromiseCentury2011}. These systems are designed as standalone receptacles where product is placed momentarily during examination or as probes and sensors inserted into the product or mounted on the storage containers for continuous monitoring. The sensors require minimal training to operate and although scalable, relatively low-cost, and non-destructive, the modality is highly susceptible to external environmental noise and vibration. 

Algorithms for vibro-acoustic stored product insect detection leverage signal processing and statistical modeling to utilize the characteristics of the audio signals generated by insects to determine the presence and severity of the infestation. These algorithms often focus on identifying specific features in the acoustic signals, such as frequency content, impulse patterns, and other spectral and temporal characteristics that are indicative of insect activity. More recently, emerging machine learning and artificial intelligence capabilities have been applied to recognize insect audial signatures in raw data, spectrograms, and other transformations. Development of these algorithms is challenged by the need for large amounts of diverse and annotated data as well as the variability in insect sounds across species, life stages, infested materials, and environmental conditions. Vibro-acoustic insect detection systems often employ uncalibrated sensors, which complicates a reliable comparison of signatures. Operational insect detection within warehouse and inspection facilities is frequently confronted by elevated ambient noise levels, which contribute to increased false positive and negative rates. Consequently, robust methods to assess and mitigate their influence on measurement accuracy must be developed. 

This thesis presents the  Acoustic Stored Product Insect Detection System (A-SPIDS) which was developed to address the need for an efficient alternative to manual visual inspection of stored products for insects. Originally funded by the Department of Homeland Security (DHS) Cross-Border Threat Screening and Supply Chain Defense (CBTS) to research and prototype alternative low-cost sensors, this work expanded into deeper analysis using the recordings made by the system of live insect specimen in various materials to study their acoustic characteristics, explore normalization methods, investigate reducing the influence of background noise, and implement a detection algorithm that could detect a single insect within \SIrange{1}{1.5}{\liter} of product. 

A-SPIDS uses piezoelectric disc sensors to capture the acoustic signals generated by insects as they move around the product. The system was used to collect an extensive amount of recordings using live insect specimens at various life stages in materials of different kernel size and hardness while exposed to varying levels of natural and simulated background noise. A database package was created to assist in collecting, storing, organizing, and retrieving the recorded audio and its associated metadata for use in analysis, visualization, algorithm development, and future deployment.

\newpage
The impulse signals generated by the insects as they maneuvered around the material were analyzed to identify unique spectral and temporal characteristics. Since the piezoelectric discs are not calibrated, a method of normalization was introduced that used the self-noise of the system as a reference level to compare the respective signals. The self-noise normalization could be applied when the total energy of the insect signal is less than the energy of the sensor noise. The experimental data confirmed that the self-noise remained after the removal of the impulse signals from the recording and that it could be used as a normalization reference. It was also applied to recordings made with different sensor models from the BugBytes dataset \cite{mankinBugBytesSound2019}.

The normalization methods enabled the derivation of two metrics to analyze and compare the acoustic properties of the recorded signals. The Normalized Strong Pulse Amplitude (NSPA) and the Normalized Signal Energy Level (NSEL) use the amplitude and spectral energy components, respectively. NSPA presents the signal noise ratio (SNR) for impulses generated by insects above the system noise level and allows for comparison of detection ranges across the different insect-material pairings when using a peak pulse amplitude threshold algorithm. Since the NSPA parameter depends on a defined frequency band for filtering and normalization, it was used to identify optimal ranges to minimize the influence of background noise. The NSEL characterizes the averaged energy of the insect signal compared to the system noise level. It could characterize the SNR of the detection system based on signal energy detection. These metrics, when statistically analyzed, provided insight to the variation between the insect-material pairings and a foundation towards an effective detection algorithm. Extensive testing revealed that impulse-based algorithms, as indicated by consistently higher NSPA values compared to NSEL metrics, are more effective for insect detection.

Although A-SPIDS was built with physical insulation that was shown to reduce external noise by an average of \SI{45}{\deci\bel} above \SI{2000}{\hertz}, signals similar to those produced by insects still penetrate through to the sensitive piezoelectric sensors and would result as false positives. The developed NSPA algorithm was used to estimate the amplitudes of the false positive, insect-like signals and determine that a filter range of \SIrange{1565}{6000}{\hertz} would result in a \SI{99.4}{\percent} detection rate when exposed to noise levels up to \SI{80}{\deci\bel A} in the curated dataset.

In order to accommodate for further resilience to external noise, a noise removal algorithm was developed leveraging the microphones included in the system design. The algorithm identifies high peak-to-RMS signals outside the sensor and eliminates them from the internal piezoelectric sensor waveforms. This leads to a \SI{100}{\percent} detection rate for the available recording subset with a zero false positive rate at noise levels up to \SI{100}{\deci\bel A}.

This research provides several unique and impactful contributions to the field of acoustic sensing for insect detection. It demonstrates the feasibility of using simple signal processing techniques for effective insect detection without relying on complex machine learning models. The usage of filters and external microphones enabled the sensor to effectively capture and isolate insect sounds from background noise while maintaining a small and portable footprint, ensuring its practicality for real-world applications. It also makes publicly available a large dataset of acoustic recordings that can be expanded and utilized by other researchers in the field, especially for AI applications. 

The developed sensor and algorithms have the potential to enhance and expedite the inspection process at Ports of Entry, enabling CBPAS to process shipments more efficiently, accurately, and rapidly while ensuring protection of national agriculture and ecologies. The system can be utilized at other points of the stored product supply chain including farms, processing plants, storage facilities, transportation vehicles, and retail locations. Furthermore, the methodologies and insights derived from this research could inform the development of analogous acoustic sensing solutions for detecting invasive species in other high-risk vectors, including fruits and vegetables, wood and wood products, and live plant materials.

The thesis is presented as follows. The first chapter describes the design and development of the A-SPIDS sensor system and the live insects, materials, and simulated background noise conditions used during the recording as well as the methodology for collecting, storing, and organizing the recorded data and its associated metadata. The following chapters delve into the analysis of the acoustic characteristics of the recorded insect sounds, introducing normalization methods and metrics to compare the signals across different insect-material pairings. The subsequent chapter focuses on the development of the detection algorithm that leverages the insights gained from the previous sections to effectively identify insect presence within stored products, including the preliminary usage of NSPA as a threshold criteria, addressing the influence of background noise, and the implementation of a noise removal algorithm using the external microphones. Finally, the conclusion summarizes the key findings, discusses the implications of the research, and outlines potential directions for future work in this field.

\newpage

\end{document}